\chapter{Arquitectura Unificada QFChain-All-in-One}
\label{chap:allinone}

La fase final de la implementación tecnológica de TCCU-6D se materializa en el ecosistema \textbf{QFChain-All-in-One}. Esta arquitectura integra de forma nativa la red de validación, el puente informacional con el simulador 6D, y las interfaces de usuario (móvil y web) bajo un único marco operativo.

\section{Componentes del Ecosistema}
El sistema se organiza en cuatro pilares fundamentales:
\begin{itemize}
    \item \textbf{Core (Node):} Motor de consenso basado en la Constitución v6.5, encargado de la producción de bloques y validación de coherencia.
    \item \textbf{TCCU Bridge:} Puente de enlace en Python que sincroniza los eventos del simulador 6D con registros inmutables en la blockchain.
    \item \textbf{Explorer:} Interfaz web para la auditoría en tiempo real del campo informacional $\Phi$ y la distribución de reputación.
    \item \textbf{Wallet Mobile:} Aplicación Android para la gestión de soberanía cognitiva y participación en el Oráculo de Consenso desde dispositivos móviles.
\end{itemize}

\section{El Puente Informacional (tccu\_bridge.py)}
El puente actúa como el traductor entre el devenir del simulador y la memoria de la red. Utiliza llamadas JSON-RPC para registrar estados ontológicos.

\begin{lstlisting}[language=Python, caption=Interfaz del Puente TCCU-QFChain]
def send_transaction(from_addr, to_addr, value, data_hex):
    # Traduce el estado ontologico a una transaccion de blockchain
    params = [{
        "from": from_addr,
        "to": to_addr,
        "value": str(value),
        "data": "0x" + data_hex
    }]
    return rpc_call("qfchain_sendTransaction", params)
\end{lstlisting}

\section{Sincronización del Tiempo Cósmico}
El número de bloque en QFChain se utiliza como el \textbf{reloj de sincronización} para las dimensiones temporales $t_1, t_2, t_3$ del simulador. Esto garantiza que todos los observadores en la red perciban una evolución coherente del universo simulado.

\section{Conclusión de la Integración}
La arquitectura All-in-One cierra el ciclo de la Creación Continua Universal, proporcionando no solo el modelo teórico y la simulación, sino la infraestructura completa para una red cognitiva soberana. El universo simulado y la red de gobernanza son ahora una única entidad digital inseparable.
